\documentclass[11pt]{article}

% ---------------------------------------------------------------------------
%  L7 — The subdominant exponential and the amplitude ratio of convex
%  animals.  Split out of L5 (the squeeze, area and semiperimeter results stay
%  there); L5 and L7 are companions and each is readable without the other's
%  proofs.
%  Category L (entirely machine-written); see paper/README.md and
%  docs/publication-split.md.
%
%  N3 COLLISION.  Gouyou-Beauchamps & Leroux, FPSAC 2004, §2.3, have the block
%  decomposition and the mirror equality of \label{prop:mirror}, for convex
%  polyominoes on the honeycomb lattice.  docs/publication-split.md requires
%  that citation in three places across the pair; the block-decomposition
%  placements live in L5 and the mirror-equality placements live here, marked
%  "attribution --- k of 2": rem:attr2 (at prop:mirror) and the paragraph
%  opening \label{sec:related}.  Do not let a revision drop one of them, and do
%  not let the L5/L7 split reduce the total below three.
%
%  Source material, all in this repository:
%    results/hv-growth-sandwich.md      Props 6-11, Lemmas 1-4, Conjecture 8
%    results/convex-polyplets.md        mu, nu, the PSLQ boxes
%    results/middle-kingdom-phase3.md   the grid, Corollary 4, Proposition 5
%    results/mk-dir4-perimeter.md       the dir4 series
% ---------------------------------------------------------------------------

% ---------------------------------------------------------------------------
%  paper/shared/preamble.tex — packages and macros common to every manuscript
%  in this directory.  Factored out of paper/polyplets-report.tex and
%  paper/technical-report.tex, which between them had two copies of the same
%  twenty lines.
%
%  technical-report.tex does NOT \input this file and is not to be edited to do
%  so: it is jasonp's prose and is read-only to the machine (paper/README.md,
%  "Who may edit what").  The duplication there is deliberate.
%
%  Assumes \documentclass[11pt]{article} has already been issued.
% ---------------------------------------------------------------------------
\usepackage[margin=1in]{geometry}
\usepackage{amsmath,amssymb}
\usepackage{amsthm}
\usepackage{booktabs}
\usepackage{graphicx}
\usepackage{numprint}
\usepackage{microtype}
\usepackage[table]{xcolor}
\usepackage[hidelinks]{hyperref}
\usepackage{flafter}   % a float never appears before the text that introduces it
\usepackage{float}     % [H] pins tables/figures exactly where written
\usepackage[font=small,labelfont=bf,labelsep=period]{caption}
\setlength{\intextsep}{16pt plus 3pt minus 3pt}

\npthousandsep{,}
\npfourdigitsep
\newcommand{\num}[1]{\numprint{#1}}
\newcommand{\g}{\cellcolor{gray!15}}

% Theorem environments.  One counter shared across kinds, so that a reference
% like "Theorem 4" is unambiguous in a paper that also has lemmas and
% propositions.
\theoremstyle{plain}
\newtheorem{theorem}{Theorem}
\newtheorem{proposition}[theorem]{Proposition}
\newtheorem{lemma}[theorem]{Lemma}
\newtheorem{corollary}[theorem]{Corollary}
\theoremstyle{definition}
\newtheorem{definition}[theorem]{Definition}
\newtheorem{example}[theorem]{Example}
\theoremstyle{remark}
\newtheorem{remark}[theorem]{Remark}
\newtheorem{conjecture}[theorem]{Conjecture}
\newtheorem{openproblem}[theorem]{Open Problem}

% Repo-path citations.  Every one of these must resolve in the tree that ships
% with the paper; `make gate-citations` checks the markdown, and
% paper/verify_claims.py checks the manuscripts.
\newcommand{\repo}[1]{\texttt{#1}}

% OEIS links.
\newcommand{\oeis}[1]{\href{https://oeis.org/#1}{#1}}

% Growth constants, so a rename is one edit.
\newcommand{\lam}{\lambda}
\newcommand{\muH}{\mu_H}

% ---------------------------------------------------------------------------
%  paper/shared/disclosure.tex — the two authorship disclosure blocks.
%
%  docs/publication-split.md §1 fixes the rule these implement:
%
%    "Under no circumstances will there be any ambiguity about whether a paper
%     is wholly my work (none), merely my prose (some) or entirely LLM-authored
%     (some)."   — jasonp, 2026-08-07
%
%  There are exactly two blocks and they live in exactly one file, so that a
%  paper cannot quietly soften its own disclosure.  The fixed sentences take no
%  arguments.  The ONE thing a paper supplies is the per-result verification
%  ledger, because that is the part that differs honestly between papers — and
%  §1 requires it to be per result, not in aggregate.
%
%  Do not add a \Pdisclosure or \Ldisclosure variant with an optional softening
%  argument.  If a paper does not fit one of these two blocks, the paper is
%  wrong, not the block.
% ---------------------------------------------------------------------------

\newsavebox{\disclosurebox}

\newenvironment{disclosureframe}
  {\begin{center}\begin{lrbox}{\disclosurebox}\begin{minipage}{0.93\textwidth}\small}
  {\end{minipage}\end{lrbox}\fbox{\usebox{\disclosurebox}}\end{center}}

% --- P: jasonp's prose -------------------------------------------------------
% Byline: Jason Parker, alone.  Argument: the per-result ledger of what the
% machine did.
\newcommand{\Pdisclosure}[1]{%
\begin{disclosureframe}
\textbf{Authorship.}\enspace Every sentence of this paper was written by its
named author. He has read every proof it states and believes each one, and he
is responsible for the correctness of what it claims.

\smallskip
\textbf{Machine contribution.}\enspace #1
\end{disclosureframe}}

% --- L: entirely machine-written --------------------------------------------
% Argument: the per-result verification ledger — for each result, what warrants
% it (Lean kernel, exact certificate, or labelled measurement) and how much of
% it a human has checked.
\newcommand{\Ldisclosure}[1]{%
\begin{disclosureframe}
\textbf{Authorship.}\enspace This document was written end to end by a large
language model. That includes its mathematics, its prose, its tables and its
\LaTeX{}. No sentence of it was written by a human. The mathematics it reports
was likewise derived by machine.

\smallskip
\textbf{Human role.}\enspace The person who directed this work chose what to
compute, chose what to publish, and is responsible for the decision to release
this document. Except where the ledger below says otherwise, that person has
not personally checked the arguments here, and this disclosure is not to be
read as an endorsement of them.

\smallskip
\textbf{What warrants each result.}\enspace #1
\end{disclosureframe}}

% --- Draft banner ------------------------------------------------------------
% For an L-paper whose verification ledger is not yet settled.  Loud on purpose:
% an L-paper without a truthful ledger must not be mistaken for a finished one.
%
% \firstdraftbanner is the standard wording, which five papers previously
% carried character-for-character in their own preambles.  It lives here for the
% same reason the disclosure blocks do: identical text in seven places drifts,
% and a paper must not be able to soften it by editing its own copy.  The
% optional argument is for a gate specific to one paper, appended verbatim; a
% paper needing different wording (L6, compute-gated) calls \draftbanner direct.
\newcommand{\firstdraftbanner}[1][]{%
\draftbanner{This is a first draft. The verification ledger below records what
warrants each result \emph{mechanically}; the column that records what a human
has read is empty, deliberately, because nobody has read it yet. Do not
circulate until that column says something true.#1}}

\newcommand{\draftbanner}[1]{%
\begin{center}
\fcolorbox{red!60!black}{yellow!15}{%
  \begin{minipage}{0.93\textwidth}\small
  \textbf{DRAFT — NOT FOR CIRCULATION.}\enspace #1
  \end{minipage}}
\end{center}}


\newcommand{\Z}{\mathbb{Z}}
\newcommand{\Q}{\mathbb{Q}}
\newcommand{\R}{\mathbb{R}}
\newcommand{\HV}{\mathrm{HV}}

\title{\textbf{The subdominant exponential and the amplitude ratio of convex
polyplets}\\[4pt]
\large A second growth constant, and a closed expression for a ratio
previously known only as digits}
\author{[byline: jasonp's call --- \repo{docs/publication-split.md} \S1]}
\date{August 2026}

\begin{document}
\maketitle

\firstdraftbanner[ One further gate is specific to this paper and its
companion: the attribution to Gouyou-Beauchamps and Leroux
(\S\ref{sec:related}) must survive every revision, in both of the places it
appears here.]

\Ldisclosure{%
\emph{Lemma~\ref{lem:T} and Proposition~\ref{prop:split} (the exact split, and
both rates)} --- proved below, from Fekete's lemma and the companion paper's
three-block factorisation. Every step of Lemma~\ref{lem:T} was additionally
checked against the enumerated block through $n = 16$, against two independent
oracles, with the slack in each step measured and reported: steps (i) and (iv)
run about a factor of $2$ from failing and step (vi) a factor of $34$, so the
check confirms the shape of each step and not its constants
(\repo{results/l-paper-proof-audit.md}).
\emph{Propositions~\ref{prop:mirror} and \ref{prop:phi}} --- proved below.
\emph{Proposition~\ref{prop:ratio} (the amplitude ratio)} --- proved,
conditional on the two amplitudes existing, which is strictly weaker than
Conjecture~\ref{conj:sharp}.
\emph{Conjecture~\ref{conj:sharp}} --- \textbf{a conjecture}, and the paper is
explicit that the repository has no proof of a sharp asymptotic for \emph{any}
series in this family.
\emph{$\nu$, $r$, the Prony spectra, the amplitudes} --- labelled measurement,
with the trusted digit count itself measured rather than asserted.
\emph{The PSLQ verdicts} --- labelled measurement, stated as falsifiable box
exclusions and never as ``$r$ is transcendental''.
\emph{Inherited from the companion paper} --- the squeeze, the stack bound and
the three-block factorisation are proved there and used here as stated.
\emph{Novelty} --- the mirror equality is \textbf{a real collision} with
\cite{gblRoux2004}, cited in two places here.
\emph{Human verification:} none, as of this draft.}

\begin{abstract}
\noindent
A \emph{convex polyplet} is a king-connected set of cells of $\Z^2$ in which
every row and every column is a single contiguous run. Its companion paper
proves that every class between the staircase polyplets and the HV-convex ones
grows at the same constant $\mu = 3.128943269730886\dots$ This paper is about
what sits underneath that leading term.

The four-cone-directed subclass splits exactly into an ascending half growing at
$\mu$ and a descending half growing at a second constant
$\nu = 2.5145796\dots$, both rates proved. Read by Prony's method, the directed
class's exponential spectrum is the two halves' spectra interleaved, and the
interleaving continues past the fourth root.

The amplitude ratio of the directed and unrestricted classes, previously known
only as $54$ measured digits, has a closed expression: a ratio of two explicit
vectors against $\varphi$, the bounded solution of the three-term recurrence
$\varphi(h) = (2 - x^{h-1})\varphi(h{-}1) - \varphi(h{-}2)$. The same recurrence
computes $\mu$ by shooting, reaching $987$ digits in $1.6$ seconds with no
series and no extrapolation, against a $700$-term enumeration plus
extrapolation for $200$. We report $r$ to $251$ digits and find no integer
relation for it in any box the precision supports.

The measured spectrum turns out not to be a list of decimals: all four
exponentials, the negative one included, are reciprocals of zeros of the
$q$-Bessel-type kernel that solves the area-refined functional equation, which
makes the subdominant ratio a ratio of two kernel zeros.

What none of this reaches is a sharp asymptotic. The repository still has no
proof that any series in this family satisfies $A(n) \sim C\mu^n$; what the
kernel changes is that the missing step is now named, simplicity of the
remaining poles and control of the tail, rather than absent.
\end{abstract}

\tableofcontents

\section{Introduction}

Convex polyplets come in a range of classes: the staircase polyplets at the
bottom, the HV-convex ones at the top, and directed conditions of various
strengths in between. Companion paper L5, \emph{Convex polyplets}, proves
that all of them grow at one constant, $\mu = 3.128943269730886\dots$, so the
leading term distinguishes none of them. Everything that does distinguish them
sits below that term, and this paper is about what sits there.

Two results. The four-cone-directed class splits exactly in two, and the halves
grow at different rates: the ascending half at $\mu$, the descending half at a
second constant $\nu = 2.5145796\dots$ Both rates are proved
(\S\ref{sec:subdominant}), and Prony's method reads the whole exponential
spectrum of each series, showing the directed class's spectrum to be the two
halves' interleaved. Second, the amplitude ratio of the directed and
unrestricted classes, which the repository had carried for a long time as $54$
measured digits and nothing else, has a closed expression
(\S\ref{sec:amplitude}): a ratio of two explicit vectors against the bounded
solution of a three-term recurrence. That recurrence also computes $\mu$ itself
by shooting, faster and further than the enumeration it replaces.

What none of it reaches is a sharp asymptotic, and \S\ref{sec:notbuy} is
explicit that the gap is not peculiar to the subdominant: no series in this
family is known to satisfy $A(n) \sim C\mu^n$ at all.

\section{The class and its blocks}
\label{sec:setup}

\begin{definition}
A \emph{convex polyplet}, equivalently an \emph{HV-convex polyplet} or convex
animal on the king lattice, is a finite king-connected set of cells of $\Z^2$ in
which every row and every column is a single contiguous run, counted up to
translation. Write $A(n)$ for the number with $n$ cells. A \emph{staircase}
polyplet, the king-lattice analogue of the parallelogram polyominoes, is
one in which, scanning columns left to right, both the bottom and the top of
the column interval are nondecreasing; write $M(n)$ for their number
(OEIS \oeis{A225114}).
\end{definition}

An HV-convex polyplet is a sequence of column intervals $[b(j), t(j)]$,
$j = 1,\dots,k$, with $b$ valley-unimodal and $t$ peak-unimodal, counted up to
translation. Write
\[
  d(j) = b(j{+}1) - b(j), \qquad h(j) = t(j) - b(j) + 1, \qquad
  s(j) = h(j) - h(j{+}1).
\]
Give column $j$ the \emph{phase} $(p_b, p_t)$, where $p_b = 1$ if some earlier
step had $d > 0$ and $p_t = 1$ if some earlier step had $t$ decreasing.
Unimodality makes both bits monotone in $j$, so the transfer operator is
block-triangular in the four phases:

\begin{table}[H]
\centering\small
\begin{tabular}{c l l l}
\toprule
phase & meaning & step constraint & kernel $K(h,h')$ \\
\midrule
$(0,0)$ & $b$ falls, $t$ rises & $s \le d \le 0$ & $h'-h+1$, \; $h' \ge h$ \\
$(1,0)$ & both rise & $d \ge \max(0,s)$ & $\min(h,h')+1$ \\
$(0,1)$ & both fall & $d \le \min(0,s)$ & $\min(h,h')+1$ \\
$(1,1)$ & $b$ rises, $t$ falls & $0 \le d \le s$ & $h-h'+1$, \; $h' \le h$ \\
\bottomrule
\end{tabular}
\caption{The four diagonal blocks. Blocks $(1,0)$ and $(0,1)$ are the staircase
kernel and its mirror; blocks $(0,0)$ and $(1,1)$ are stacks.}
\label{tab:blocks}
\end{table}

The \emph{four-cone-directed} (dir4) condition restricts the step set further;
its effect on the table is confined to the $(0,0)$ and $(0,1)$ blocks, which is
the fact \S\ref{sec:subdominant} starts from.

Three results of companion paper L5, \emph{Convex polyplets}, are used
without reproof:

\begin{itemize}
\item \textbf{The squeeze.} Every class $\mathcal C$ with
  $\text{staircase} \subseteq \mathcal C \subseteq \text{HV-convex}$ has
  $\lim_n \mathcal C(n)^{1/n} = \mu = 3.128943269730886\dots$
\item \textbf{The stack bound.} The number $P(n)$ of monotone-height blocks of
  area $n$ satisfies $P(n) \le E(n) := (n{+}1)^{4\sqrt n + 6}$, so
  $P(n)^{1/n} \to 1$.
\item \textbf{The three-block factorisation.}
  $A(n) \le 2(n{+}1)^2\sum_{i+j+l=n}P(i)M(j)P(l)$, from cutting the column
  sequence at its two phase transitions.
\end{itemize}

\section{The subdominant structure}
\label{sec:subdominant}

The squeeze settles the leading term and says nothing about the corrections,
which is where the four-cone class differs from the unrestricted one.
Table~\ref{tab:blocks} locates the corrections too: of the four diagonal blocks,
exactly two carry exponential weight standalone.

\begin{table}[H]
\centering\small
\begin{tabular}{c l l}
\toprule
block & unrestricted & under the four-cone condition \\
\midrule
$(0,0)$ & sub-exponential (L5, stack bound) & still sub-exponential \\
$(1,0)$ & staircase, growth $\mu = 3.1289\dots$ & untouched \\
$(0,1)$ & mirror staircase, growth $\mu$ & \textbf{truncated, growth $\nu = 2.5146\dots$} \\
$(1,1)$ & sub-exponential (L5, stack bound) & untouched \\
\bottomrule
\end{tabular}
\caption{The truncated $(0,1)$ block is the exception, and the only one.}
\label{tab:subdom}
\end{table}

The truncated $(0,0)$ block stays sub-exponential for a \emph{different} reason
from L5's stack bound: with $h' \ge h$, the entry count
$\min(2, h'-h+1)$ is $2$ only at a strict rise, strict rises take distinct
heights, so there are at most $\sqrt{2n}$ of them and the weight is at most
$2^{\sqrt{2n}}$. The truncated $(0,1)$ block is not height-monotone ($h' \le
h$ at weight $2$, $h' = h+1$ at weight $1$, $h' > h+1$ forbidden), so heights
can climb one row at a time and the stack bound does not apply. Counted
by area it is
\[
  T(n) \;=\; 1,\, 3,\, 8,\, 21,\, 54,\, 138,\, 350,\, 885,\, 2233,\, 5626,\, \dots
\]
with no OEIS match on ten terms. (Its six-term prefix collides with
\oeis{A127358} and two others, all of which diverge by $n = 9$: $2230$ against
$2233$.)

\subsection{The two rates}

Partition the four-cone class by whether an animal's phase path ever visits
$(0,1)$:
\[
  D(n) \;=\; D_{\mathrm{asc}}(n) + D_{\mathrm{desc}}(n).
\]

\begin{lemma}
\label{lem:T}
$\lim_n T(n)^{1/n}$ exists. Write $\nu$ for it.
\end{lemma}

\begin{proof}
Let $T_h(n)$ count runs whose first column has height $h$, and write
$T_1 = T_{h=1}$.
(i) $T(n) \ge 2T(n{-}1)$: append a height-$1$ column, which the truncated kernel
admits at weight $2$ from any height.
(ii) $T_h(n) \le 2T(n{-}h)$: delete the first column, the step out of it having
weight at most $2$. With (i), $T_h(n) \le 2^{2-h}T(n{-}1)$, so
$\sum_{h\ge4}T_h(n) \le T(n{-}1)/2 \le T(n)/4$ and
$T(n) \le 4\max_{h\le3}T_h(n)$.
(iii) $T_h(n) \le T_1\bigl(n + h(h{-}1)/2\bigr)$: prepend the climb
$1,2,\dots,h{-}1$, every step of which is the weight-$1$ transition $h' = h+1$;
for $h \le 3$ that costs at most $3$. So $T(n) \le 4T_1(n{+}3)$.
(iv) $T_1(i)T_1(j) \le T_1(i{+}j)$: join two runs that both start at height $1$
by the step ``last height $\to 1$'', weight $2$; the pair is recovered from the
join and the prefix area.
Fekete gives $\lim T_1(n)^{1/n} = \sup_n T_1(n)^{1/n}$, and
$T_1 \le T \le 4T_1(\cdot{+}3)$ squeezes $T$ onto the same limit.
\end{proof}

\begin{proposition}
\label{prop:split}
$\lim_n D_{\mathrm{asc}}(n)^{1/n} = \mu$ and
$\lim_n D_{\mathrm{desc}}(n)^{1/n} = \nu$.
\end{proposition}

\begin{proof}
$D_{\mathrm{asc}} \ge M$ termwise: a staircase polyplet has $d \ge 0$ and
$d \ge s$ at every step, so it can never enter $(0,1)$, which needs $d \le 0$
and $d < s$. With $D_{\mathrm{asc}} \le D$ and L5's squeeze, the first claim follows.

For the second, L5's three-block factorisation with $T$ in place of $M$
gives $D_{\mathrm{desc}}(n) \le 2(n{+}1)^4E(n)^2C_\varepsilon(\nu+\varepsilon)^n$,
so $\limsup D_{\mathrm{desc}}(n)^{1/n} \le \nu$. Conversely take a $T$-run of
area $m$ whose first column is $[0,h{-}1]$ and prepend the column $[0,h]$: the
joining step has $d = 0 \le 0$ and $s = 1 > d$, so it enters $(0,1)$, the run's
own steps keep it there, and the result is a (dir4, HV-convex) animal of area
$m + h + 1$. Restricting to $h = 1$ gives $D_{\mathrm{desc}}(n) \ge T_1(n{-}2)$,
so $\liminf D_{\mathrm{desc}}(n)^{1/n} \ge \nu$ by Lemma~\ref{lem:T}.
\end{proof}

Measured,
\begin{quote}\small
$\nu = 2.5145796438787291885104371943430998201410308559007832719354957107238\\
4099229632690389299337761749953732175539964789209170137762869374741150849461805388232$
\end{quote}
to $153$ trusted digits by the conservative extrapolation and $242$ by Prony's
method on the same terms.

\subsection{The Prony spectrum}

Reading the exponential spectrum of each $700$-term series by Prony's method:
solve for the constant-coefficient recurrence the tail obeys best and take the
roots of its characteristic polynomial, so every $\lambda_i$ is measured in one
solve, with no peeling and no assumed model:

\begin{table}[H]
\centering\small
\begin{tabular}{l l l l l}
\toprule
series & $\lambda_1$ & $\lambda_2$ & $\lambda_3$ & $\lambda_4$ \\
\midrule
staircase & $3.12894326973088\dots$ & $1.50504922775900\dots$ & $1.28433727098118\dots$ & $-1.25776216033063\dots$ \\
HV-convex & same & same & same & same \\
$D_{\mathrm{asc}}$ & same & same & same & same \\
block $T$ & $2.51457964387872\dots$ & $1.43040460381247\dots$ & $1.24846593371011\dots$ & $1.17441805313317\dots$ \\
$D_{\mathrm{desc}}$ & same & same & same & same \\
$D$ = dir4 & $3.12894326973088\dots$ & $2.51457964387872\dots$ & $1.50504922775900\dots$ & $1.43040460381247\dots$ \\
\bottomrule
\end{tabular}
\caption{$D$'s spectrum is the other two interleaved, and the interleaving
continues past $\lambda_4$. Agreements run to $310$, $246$ and $236$ digits on
the leading claims, each against the weaker of the two trusted counts.}
\label{tab:prony}
\end{table}

Two negative controls are required to fail, and do: the block's $\lambda_1$ and
$\lambda_2$ match no root of $D_{\mathrm{asc}}$, to $0$ and $1$ digits, and the
staircase's $\lambda_2$ matches no root of $D_{\mathrm{desc}}$, to $1$.

\subsection{The spectrum is the kernel's zero set}
\label{sec:kernelspectrum}

Prony's method measures those exponentials without any model. There is a model,
and it accounts for all four.

The companion paper's \S``The kernel, and what it certifies'' solves the
area-refined functional equation by Temperley's method and reduces the
generating function to $c_0 + c_1\alpha/K$, meromorphic on the unit disc with
poles only at zeros of $K$ and of a $2\times2$ determinant. The leading
exponential is $1/q_1$ with $q_1$ the smallest positive zero of $K$. The rest
of the spectrum is the rest of the zeros:

\begin{table}[H]
\centering\small
\begin{tabular}{l l l}
\toprule
measured by Prony & zero of $K$ & $1/q$ \\
\midrule
$\lambda_1 = 3.12894326973088\dots$ & $q_1 = 0.3195967180593874655$ & $3.128943269730886252$ \\
$\lambda_2 = 1.50504922775900\dots$ & $q_2 = 0.6644300940833578092$ & $1.505049227759003935$ \\
$\lambda_3 = 1.28433727098118\dots$ & $q_3 = 0.7786116798090276394$ & $1.284337270981181429$ \\
$\lambda_4 = -1.25776216033063\dots$ & $q_4 = -0.7950628755893912422$ & $-1.257762160330635483$ \\
\bottomrule
\end{tabular}
\caption{Agreement to every digit Table~\ref{tab:prony} prints, and to $39$ on
$\lambda_1$. $\lambda_4$ is negative, so it can only come from a negative zero;
the certification work scanned the positive axis only and never saw it.}
\label{tab:kernelzeros}
\end{table}

Two consequences. The subdominant ratio $\rho = \lambda_2/\lambda_1$ of
Table~\ref{tab:tableB} is $q_1/q_2$, a ratio of two kernel zeros rather than a
measured decimal. And the sign in $\lambda_4$, which a coefficient-level reading
can only record, is the sign of a zero.

This does \emph{not} prove Conjecture~\ref{conj:sharp}. A
meromorphic function with poles at $q_1, q_2, \dots$ gives the stated expansion
once one knows the poles are simple, that nothing else contributes, and that the
tail is controlled; the first is certified at $q_1$ and the rest is the analytic
bookkeeping L5 states as an open problem. The conjecture's \emph{shape} is now
explained; its status is unchanged.

\subsection{Amplitudes and error terms}
\label{sec:notbuy}

\begin{conjecture}
\label{conj:sharp}
$D_{\mathrm{asc}}(n) = C\mu^n(1 + O(\rho^n))$ with $\rho = 0.48100879370959\dots$
and $D_{\mathrm{desc}}(n) = C'\nu^n(1 + O(\sigma^n))$ with
$\sigma = 0.568844421888\dots$, where $C = 0.4503031857123411823536\dots$ and
$C' = 2.5795006951239199769\dots$
\end{conjecture}

Granting it, $D(n) = C\mu^n + C'\nu^n + O((\mu\rho)^n + (\nu\sigma)^n)$ with
$\mu\rho = 1.50504\dots$ and $\nu\sigma = 1.43040\dots$, both below $\nu$, so the
subdominant exponential of $D$ is exactly $\nu$.

The gap is not peculiar to the subdominant. The repository has no proof that
\emph{any} series in this family has a sharp asymptotic at all:
L5's squeeze and staircase supermultiplicativity give exponential
rates, never $A(n) \sim C\mu^n$, and L5's tabulated amplitudes are all measured, none derived. Asking for the second
exponential as a theorem is asking for strictly more than is known about the
first.

A route exists and it is specific. The \oeis{A225114} entry carries a
\emph{conjectured} continued fraction (Kurkov, September 2024),
\[
  \text{g.f.} = \cfrac{1}{2 - \cfrac{1}{1 - \cfrac{x}{1 - \cfrac{x}{1 - \cfrac{x^2}{1 - \cfrac{x^2}{1 - \cdots}}}}}},
\]
whose depth-$24$ convergent reproduces $41$ terms and whose poles, as $1/x$,
land on the measured staircase spectrum, to $30$, $27$, $14$ and $16$ digits on
$\lambda_1,\dots,\lambda_4$, improving with depth. A proof of that continued
fraction would supply the meromorphic continuation and the spectral gap that a
coefficient-level argument cannot reach.

\subsection{What the $d_n/d_{n-1}$ diagnostic reports}

The two series share their subdominant correction; the four-cone series carries
an extra exponential in front of it, which the diagnostic reports instead:

\begin{table}[H]
\centering\small
\begin{tabular}{l l l}
\toprule
series & $d_n/d_{n-1}$ & what it is \\
\midrule
HV-convex, staircase & $0.481008794$ & $\lambda_2/\lambda_1$, the staircase block's own correction \\
$D_{\mathrm{asc}}$ & $\mathbf{0.48100879371}$ & the same one \\
$D_{\mathrm{desc}}$, block $T$ & $0.568844421888$ & the block's own correction \\
(dir4, HV-convex) & $0.803651401483$ & $\nu/\mu$, the \emph{extra} exponential in front \\
\bottomrule
\end{tabular}
\caption{The restricted series has an extra exponential in front of the shared
one, and the $d_n/d_{n-1}$ diagnostic reports only the largest correction
present.}
\label{tab:tableB}
\end{table}

The second row is measured on four-cone-directed animals directly, with no
peeling and no annihilator. Deleting the descending half stops the $2.5146$ from
polluting the extrapolation, and the same run's $\mu$ rises from $51$ trusted
digits to $200$.

This rules out a bijection between the two classes: a bijection would have to
account for a half of one side, growing at $\nu$, with no counterpart on the
other.

\section{The amplitude ratio}
\label{sec:amplitude}

The ratio $r = C_{\text{dir4}}/C_{\HV}$ was for a long time just a measured
decimal. It has a closed expression.

\begin{proposition}[the mirror halves]
\label{prop:mirror}
Partition the HV-convex animals by which middle phase their path visits. Then
$A_{(1,0)}(n) = A_{(0,1)}(n)$ for every $n$, and the remainder, the paths
$(0,0)\to(1,1)$, is sub-exponential. Hence if $C_{\HV} = \lim A(n)/\mu^n$
exists, $C_{\HV} = 2\,C(A_{(1,0)})$.
\end{proposition}

\begin{proof}
The vertical mirror $[b(j),t(j)] \mapsto [-t(j),-b(j)]$ is an area-preserving
involution of the class that leaves every column height alone. It carries
$d(j) = b(j{+}1)-b(j)$ to $-(t(j{+}1)-t(j))$, so ``$b$ has risen'' becomes
``$t$ has fallen'': the phase bits swap. Phases $(0,0)$ and $(1,1)$ are fixed
and $(1,0)\leftrightarrow(0,1)$, so the involution is a bijection between the
first two parts. The remainder factors as a phase-$(0,0)$ block joined to a
phase-$(1,1)$ block at one step, both stacks by L5's stack bound, so it is
at most $(n{+}1)^2E(n)^2 = \mu^{o(n)}$.
\end{proof}

\begin{remark}[attribution --- 1 of 2]
\label{rem:attr2}
The mirror equality of the two middle blocks is Gouyou-Beauchamps and Leroux's
$H_{02} = Pa = H_{20}$ in the convex-polyomino setting~\cite{gblRoux2004}; see
\S\ref{sec:related}. What this proposition adds here is the factor
$\tfrac12$ and the sub-exponential bound on the remainder.
\end{remark}

\begin{proposition}[the staircase eigenvector is a three-term recurrence]
\label{prop:phi}
Let $T(x)_{h,h'} = x^{h'}(\min(h,h')+1)$ for $h,h' \ge 1$ be the staircase
block's operator, and let $\varphi$ solve
\[
  \varphi(0) = 1, \quad \varphi(1) = 2, \quad
  \varphi(h) = \bigl(2 - x^{h-1}\bigr)\varphi(h{-}1) - \varphi(h{-}2)
  \quad (h \ge 2).
\]
For $0 < x < 1$: $T(x)\varphi = \varphi$ if and only if
$\varphi(h) - \varphi(h{-}1) \to 0$; that happens at exactly one $x = x_c$; and
$x_c = 1/\mu$.
\end{proposition}

\begin{proof}
Write $Q(h) = \sum_{h'>h}x^{h'}\varphi(h')$. Splitting $T(x)\varphi = \varphi$
at $h' = h$ gives $\varphi(h) = \sum_{h'\le h}x^{h'}(h'{+}1)\varphi(h') +
(h{+}1)Q(h)$; differencing once gives
$\varphi(h)-\varphi(h{-}1) = x^h\varphi(h) + Q(h)$, and differencing again gives
the recurrence. Conversely, \emph{define}
$Q(h) := \varphi(h)-\varphi(h{-}1)-x^h\varphi(h)$ from the recurrence's
solution; the second difference makes $Q(h{-}1)-Q(h) = x^h\varphi(h)$ an
identity, so $Q(h) = Q(\infty) + \sum_{h'>h}x^{h'}\varphi(h')$ and the
eigenvector equation holds exactly when $Q(\infty) = 0$. Since $x < 1$ the
recurrence's coefficient tends to $2$ and its solution space is spanned by one
asymptotically constant and one asymptotically linear solution; $Q(\infty) = 0$
says the linear one is absent, one analytic condition on $x$.

At such an $x$, $\varphi$ is positive: the differences
$\delta(h) = \varphi(h)-\varphi(h{-}1)$ obey
$\delta(h) = \delta(h{-}1) - x^{h-1}\varphi(h{-}1)$ with $\delta(1) = 1$, so
$\delta$ strictly decreases while $\varphi > 0$, and $\delta(h) \to 0$ forces
$\delta > 0$ throughout: $\varphi$ increases from $2$ to $2.5374225302\dots$
$T(x)$ is positive and compact on $\{f : |f(h)| \le C(1+h)\}$, so by
Krein--Rutman~\cite{kreinRutman1948} a positive eigenvector's eigenvalue is the
spectral radius. The staircase generating function is
$M(x) = v_0(x)(I-T(x))^{-1}\mathbf 1$ with $v_0(h) = x^h$, of radius $1/\mu$,
and the spectral radius of $T(x)$ increases continuously in $x$, so it reaches
$1$ at $x = 1/\mu$.
\end{proof}

Shooting on that recurrence computes $\mu$ to arbitrary precision in
$O(h_{\max})$ operations, with no series and no extrapolation: it reproduces
all $199$ banked digits, continues past them, and reaches $987$ digits in
$1.6$ seconds on one core, against a $700$-term enumeration plus extrapolation
for $200$.

\begin{proposition}[the ratio]
\label{prop:ratio}
Write $U(h,h') = h'-h+1$ and $U_4(h,h') = \min(2, h'-h+1)$ for $h' \ge h$ (the
$(0,0)$ block, unrestricted and four-cone-truncated), $L(h,h') = \min(h,h'{+}1)$
(the map $(0,0)\to(1,0)$, and also $(0,0)\to(0,1)$, the same matrix, which is
Proposition~\ref{prop:mirror} at the operator level), each carrying $x^{h'}$,
and
\[
  w = v_0(I-U)^{-1}L, \qquad w_4 = v_0(I-U_4)^{-1}L,
  \qquad\text{evaluated at } x_c .
\]
If $C(D_{\mathrm{asc}}) = \lim D_{\mathrm{asc}}(n)/\mu^n$ and
$C_{\HV} = \lim A(n)/\mu^n$ exist, then
\begin{equation}
  r \;=\; \frac{C_{\text{dir4}}}{C_{\HV}}
    \;=\; \frac12\cdot\frac{w_4\cdot\varphi}{w\cdot\varphi} .
  \label{eq:ratio}
\end{equation}
\end{proposition}

\begin{proof}
The phase automaton makes the path decomposition exact as generating functions:
$A_{(1,0)}(x) = w(x)(I-T(x))^{-1}q(x)$ with
$q = 1 + L_{(1,0)\to(1,1)}(I-T_{(1,1)})^{-1}\mathbf 1$, and
$D_{(1,0)}(x) = w_4(x)(I-T(x))^{-1}q(x)$ with the \emph{same} $T$ and the
\emph{same} $q$: the four-cone floor bites only inside the $(0,0)$ block,
since every step out of $(1,0)$ and every step within $(1,1)$ already has
$d \ge 0$. The $(0,0)$ and $(1,1)$ blocks are stacks
(L5, stack bound), so $w$, $w_4$ and $q$ are finite at $x_c$. By
Proposition~\ref{prop:phi} and Krein--Rutman the pole of $(I-T(x))^{-1}$ at
$x_c$ is simple with rank-one residue $\varphi\psi^{\mathsf T}/\langle\psi,
\varphi\rangle$, so the two limits
$\lim_{x\to x_c^-}(1-\mu x)A_{(1,0)}(x)$ and
$\lim_{x\to x_c^-}(1-\mu x)D_{(1,0)}(x)$ are $N(w\cdot\varphi)$ and
$N(w_4\cdot\varphi)$ with the same $N$, which cancels. If the amplitudes exist
then Abel's theorem identifies those limits with $C(A_{(1,0)})$ and
$C(D_{(1,0)})$. Finally $C_{\HV} = 2C(A_{(1,0)})$
(Proposition~\ref{prop:mirror}) and $C_{\text{dir4}} = C(D_{(1,0)})$, the
$D_{\mathrm{desc}}$ and ``via neither'' parts being $o(\mu^n)$ by
Proposition~\ref{prop:split} and L5's stack bound.
\end{proof}

The hypothesis is strictly weaker than Conjecture~\ref{conj:sharp}, which asks
for the error terms as well. Equation~\eqref{eq:ratio} settles what a closed form would have to look like: a ratio of two $q$-series in $q = 1/\mu$, both computed
by $O(h_{\max})$ prefix-sum recurrences, with
$w\cdot\varphi = 1.29770192341040021939895\dots$ and
$w_4\cdot\varphi = 1.19935960397018718067118\dots$

\subsection{Three determinations of $r$}

\begin{quote}\small
$r$, $251$ trusted digits:\\[2pt]
\texttt{0.4621090492099424400356623877003058328411634397299804247065092921445}\\
\texttt{22950511444483541019804848618234062089755257041466533204384387638170139}\\
\texttt{068144091014356289934131553812178862045404673728300959682539208197483810}\\
\texttt{09549588646431928574874902351771539511812}
\end{quote}

\begin{table}[H]
\centering\small
\begin{tabular}{l l r}
\toprule
route & own cross-check & trusted \\
\midrule
$D/A$ extrapolated & raw vs Aitken, $n=600$ vs $700$ & 54 \\
$D_{\mathrm{asc}}/A$ extrapolated & raw vs Aitken, $n=600$ vs $700$ & 185 \\
identity~\eqref{eq:ratio} & $h_{\max}$ $1200$ vs $1000$ & 492 \\
\bottomrule
\end{tabular}
\caption{Three determinations of $r$.}
\label{tab:ratio}
\end{table}

The extrapolation and the identity agree to $293$ digits, and the two routes
share no arithmetic: the identity comes from an exact operator factorisation,
the extrapolation touches no operator at all. The first row reproduces the
historical $54$-digit figure character for character, which makes it the control
on the method; the $185$ of the second row is bought by the phase split.

Negative controls, all required to fail and doing so: dropping the factor
$\tfrac12$ agrees with the measurement to $0$ digits; truncating the $(0,0)$
block at $\min(3,\cdot)$ instead of $\min(2,\cdot)$ agrees to $1$; and the
$\varphi$ recurrence at $x = 1/3.13$ fails its own eigen-equation at $10^{-2}$
where $x_c$ gives $10^{-528}$.

\subsection{Is $r$ algebraic?}
\label{sec:ralg}

At $251$ digits, no integer relation in any in-capacity box:
\[
  \text{degree} \le 30 \text{ at height} \le 10^4, \;\le 20 \text{ at } 10^5,
  \;\le 10 \text{ at } 10^{11}, \;\le 5 \text{ at } 10^{20},
  \;\le 3 \text{ at } 10^{30}, \;\le 2 \text{ at } 10^{40}.
\]
Nor over $\Q(\mu)$, the natural field since $\mu$ is the constant both series
share: $\mu$-degree $\le 5$ with $r$-degree $\le 6$ at height $\le 10^2$, and
three further boxes. A positive control runs on every pass: $(1+\sqrt2)/3$
must be found at degree $2$, and is.

$\mu$'s own arithmetic status limits how much this can mean. There is no
proof that $\mu$ is irrational, let alone transcendental, so $\Q(\mu)$ is a
field of unknown degree and ``algebraic over $\Q(\mu)$'' cannot be settled
either way by any argument available here. What \eqref{eq:ratio} contributes is
a reason to expect the answer to be no: $r$ is a ratio of two $q$-series
evaluated at $q = 1/\mu$, and $\mu$ is itself the reciprocal of the zero of an
entire $q$-function (Proposition~\ref{prop:phi}), the same shape as Kurkov's
conjectured continued fraction.

\section{Relation to existing work}
\label{sec:related}

\paragraph{Attribution --- 2 of 2.} The mirror equality of
Proposition~\ref{prop:mirror} is Gouyou-Beauchamps and Leroux's
$H_{02} = Pa = H_{20}$~\cite{gblRoux2004}, \S2.3, for convex polyominoes on the
honeycomb lattice, where it sits inside the same block decomposition that
companion paper L5 attributes to them. Their square-lattice
companion~\cite{lrr1998} reaches the same objects by a different route,
so~\cite{gblRoux2004} is the source to cite.

What is not theirs, and is claimed here: the king lattice, where the
middle kernel is $\min(h,h')+1$ rather than $\min(h,h')$ and neither class is
exactly solved; the exponential split and its two proved rates
(Lemma~\ref{lem:T}, Proposition~\ref{prop:split}); the eigenvector recurrence
and the shooting method (Proposition~\ref{prop:phi}); and the amplitude-ratio
identity (Proposition~\ref{prop:ratio}).

\paragraph{The continued fraction.} The \oeis{A225114} entry carries a
conjectured continued fraction due to Kurkov (September 2024), quoted in
\S\ref{sec:notbuy}. It is the most concrete route to the meromorphic
continuation that a proof of Conjecture~\ref{conj:sharp} needs.

\section{Open problems}

\begin{openproblem}[a sharp asymptotic, for anything in this family]
Prove $A(n) \sim C\mu^n$, or the same for $M(n)$. Nothing in either paper
reaches it, and Conjecture~\ref{conj:sharp} depends on it. Two routes are now
concrete rather than one: the kernel of \S\ref{sec:kernelspectrum}, which needs
simplicity of the poles beyond $q_1$ and a controlled tail, and a proof of
Kurkov's conjectured continued fraction for \oeis{A225114}, which would supply
the meromorphic continuation and the spectral gap by a different road.
\end{openproblem}

\begin{openproblem}[does the second pole family explain $\nu$?]
\label{op:det}
The descending half's spectrum, $2.51457964$, $1.43040460$, $1.24846593$ and
$1.17441805$, appears nowhere in Table~\ref{tab:kernelzeros}. But $K$ is only
one of the two factors that can vanish: the solution's $2\times2$ determinant is
the other, and its zeros are a second pole family nobody has scanned. If $\nu$
is the reciprocal of the smallest of them, both halves of the interleaving in
Table~\ref{tab:prony} have one mechanism and this paper's central observation
becomes a statement about two zero sets rather than about two measurements.

Attempted here and \emph{not} settled: building the determinant as a truncated
$q$-series with the solver's own arithmetic and scanning it gives a smallest
positive zero that moves with the truncation order ($1.2433$ at $24$ terms,
no sign change at all at $32$, $1.1794$ at $40$ and $48$) and returns the
same value for the king and the control, which cannot be right. The
determinant's zeros need a construction that converges, not a longer scan of
this one. \repo{experiments/convex\_det\_zeros.py} is the negative attempt,
kept so the next one starts after it rather than before.
\end{openproblem}

\begin{openproblem}[the block series $T$]
$1, 3, 8, 21, 54, 138, 350, 885, 2233, 5626, \dots$ has no OEIS match and no
conjectured form. It is the truncated $(0,1)$ block, its growth constant $\nu$
is the subdominant exponential of the four-cone class, and it deserves the same
treatment \oeis{A225114} has had.
\end{openproblem}

\begin{openproblem}[is $r$ algebraic over $\Q(\mu)$?]
\S\ref{sec:ralg} excludes boxes and nothing more. Since $\mu$ is not known to be
irrational, $\Q(\mu)$ is a field of unknown degree and the question cannot be
settled either way by any argument available here.
\end{openproblem}

\section{Reproduction}
\label{sec:repro}

\begin{table}[H]
\centering\small
\begin{tabular}{l l}
\toprule
claim & check \\
\midrule
the block series $T$ & \repo{experiments/dir4\_descent\_block.py}, \repo{experiments/descent\_block\_oracle.py} \\
Table~\ref{tab:prony} & \repo{experiments/prony\_spectrum.py} \\
Propositions~\ref{prop:phi}--\ref{prop:ratio} & \repo{experiments/amplitude\_feed\_vectors.py} \\
the PSLQ boxes for $r$ & \repo{experiments/amplitude\_pslq.py} \\
$\nu$ and the amplitudes & \repo{experiments/convex\_growth.py} \\
the series this paper reads & \texttt{make build/convex\_area\_tm}; see companion paper L5 \\
\bottomrule
\end{tabular}
\caption{Where each claim is checked. Every proof-relevant item is a paper proof
above; everything in this table is a measurement or a reproduction of one.}
\label{tab:repro}
\end{table}

\bibliographystyle{plain}
\bibliography{shared/refs}

\end{document}
